% GENERATED FILE. Edit canonical lessons, then run npm run generate:books.
% canonical-source-hash: fnv1a64:fc608a7f45584831
% canonical-lessons: HI-C290-sikayat, HI-C290-riport, HI-C290-likhvana, HI-C290-daroga, HI-C290-havalat

\chapter{At the Police Station}
\label{ch:hi-at-the-police-station}
\hlchaptermodality{290}

\begin{chapteropening}
\textbf{By the end of this chapter:} \emph{I can report a crime at the police station.}

Complete the last lesson of chapter 290: I can report a crime at the police station.
\end{chapteropening}

\section[\texorpdfstring{\hi{शिकायत} (śikāyat)}{शिकायत (śikāyat)}]{\hi{शिकायत} (śikāyat) — a complaint}

\label{lesson:HI-C290-sikayat}



\begin{quote}
Before the new one: say the Hindi for on the run, then the Hindi for a chase.
\end{quote}

\subsection*{You'll want to know: \hi{शिकायत}}
\textbf{\hi{शिकायत}} — \emph{śikāyat} — \textquotedblleft{}a complaint\textquotedblright{}.

\begin{quote}
\hi{पुलिस} \hi{क्या} \hi{करेगी}, \hi{यह} \hi{आपकी} \hi{शिकायत} \hi{पर} \hi{निर्भर} \hi{करता} \hi{है।}
\end{quote}

\begin{quote}
\emph{pulis kyā karegī, yah āpkī śikāyat par nirbhar kartā hai.}
\end{quote}

\begin{quote}
What the police will do depends on your complaint.
\end{quote}

Feminine.

\begin{cousinweb}
is usually traced, through Persian, to Arabic \emph{shikāya}, \textquotedblleft{}a complaint\textquotedblright{}. \textbf{\hi{शिकायत} \hi{करना}} is to complain; at a police station, \textbf{\hi{शिकायत}} is the formal complaint.
\end{cousinweb}

\begin{grammarlens}[title={What you've built}]
One more word for telling what happened.
\end{grammarlens}

\subsection*{Your turn}
\begin{itemize}
\raggedright
  \item \emph{Say it:} \emph{śikāyat}
  \item \emph{Say it:} \emph{śikāyat}, once more
  \item \emph{Say it:} \emph{pulis kyā karegī, yah āpkī śikāyat par nirbhar kartā hai.} — What the police will do depends on your complaint.
  \item \emph{Recall:} say the Hindi for on the run, then the Hindi for a chase, then say \emph{śikāyat} again
\end{itemize}

\begin{tcolorbox}[breakable,colback=teal!4,colframe=teal!35!black,title={Before you move on}]
What does \hi{शिकायत} mean? (\textquotedblleft{}A complaint\textquotedblright{}.) Say it once more.
\end{tcolorbox}
\section[\texorpdfstring{\hi{रिपोर्ट} (riporṭ)}{रिपोर्ट (riporṭ)}]{\hi{रिपोर्ट} (riporṭ) — a report (to the police)}

\label{lesson:HI-C290-riport}



\begin{quote}
Before the new one: say the Hindi for a chase, then the Hindi for a complaint.
\end{quote}

\subsection*{You'll want to know: \hi{रिपोर्ट}}
\textbf{\hi{रिपोर्ट}} — \emph{riporṭ} — \textquotedblleft{}a report (to the police)\textquotedblright{}.

\begin{quote}
\hi{मीरा} \hi{ने} \hi{थाने} \hi{जाकर} \hi{चोरी} \hi{की} \hi{रिपोर्ट} \hi{की।}
\end{quote}

\begin{quote}
\emph{Mīrā ne thāne jākar corī kī riporṭ kī.}
\end{quote}

\begin{quote}
Mira went to the police station and reported the theft.
\end{quote}

Feminine.

\begin{cousinweb}
\textbf{\hi{रिपोर्ट}} (\emph{riporṭ}) is borrowed from English \emph{report}, taken into Hindi as a feminine noun: \textbf{\hi{चोरी} \hi{की} \hi{रिपोर्ट}}, the report of a theft, and \textbf{\hi{रिपोर्ट} \hi{करना}}, to report it.
\end{cousinweb}

\begin{grammarlens}[title={What you've built}]
One more word for telling what happened.
\end{grammarlens}

\subsection*{Your turn}
\begin{itemize}
\raggedright
  \item \emph{Say it:} \emph{riporṭ}
  \item \emph{Say it:} \emph{riporṭ}, once more
  \item \emph{Say it:} \emph{Mīrā ne thāne jākar corī kī riporṭ kī.} — Mira went to the police station and reported the theft.
  \item \emph{Recall:} say the Hindi for a chase, then the Hindi for a complaint, then say \emph{riporṭ} again
\end{itemize}

\begin{tcolorbox}[breakable,colback=teal!4,colframe=teal!35!black,title={Before you move on}]
What does \hi{रिपोर्ट} mean? (\textquotedblleft{}A report (to the police)\textquotedblright{}.) Say it once more.
\end{tcolorbox}
\section[\texorpdfstring{\hi{लिखवाना} (likhvānā)}{लिखवाना (likhvānā)}]{\hi{लिखवाना} (likhvānā) — to have written down}

\label{lesson:HI-C290-likhvana}



\begin{quote}
Before the new one: say the Hindi for a complaint, then the Hindi for a report (to the police).
\end{quote}

\subsection*{You'll want to know: \hi{लिखवाना}}
\textbf{\hi{लिखवाना}} — \emph{likhvānā} — \textquotedblleft{}to have written down\textquotedblright{}.

\begin{quote}
\hi{अरुण} \hi{ने} \hi{पुलिस} \hi{से} \hi{अपनी} \hi{शिकायत} \hi{लिखवाई।}
\end{quote}

\begin{quote}
\emph{Aruṇ ne pulis se apnī śikāyat likhvāī.}
\end{quote}

\begin{quote}
Arun had the police write down his complaint.
\end{quote}

Past, with \hi{ने} (he or she …): \hi{उसने} \hi{रिपोर्ट} \hi{लिखवाई}, \emph{usne riporṭ likhvāī}.

\begin{cousinweb}
\textbf{\hi{लिखवाना}} (\emph{likhvānā}) is the second causative of \textbf{\hi{लिखना}}, to write, which goes back to Sanskrit \emph{likhati}, \textquotedblleft{}scratches, writes\textquotedblright{} (CDIAL 11048). The ending \emph{-vā-} hands the work to someone else: \textbf{\hi{शिकायत} \hi{लिखवाई}}, got the complaint written down.
\end{cousinweb}

\begin{grammarlens}[title={What you've built}]
One more word for telling what happened.
\end{grammarlens}

\subsection*{Your turn}
\begin{itemize}
\raggedright
  \item \emph{Say it:} \emph{likhvānā}
  \item \emph{Say it:} \emph{likhvānā}, once more
  \item \emph{Say it:} \emph{Aruṇ ne pulis se apnī śikāyat likhvāī.} — Arun had the police write down his complaint.
  \item \emph{Recall:} say the Hindi for a complaint, then the Hindi for a report (to the police), then say \emph{likhvānā} again
\end{itemize}

\begin{tcolorbox}[breakable,colback=teal!4,colframe=teal!35!black,title={Before you move on}]
What does \hi{लिखवाना} mean? (\textquotedblleft{}To have written down\textquotedblright{}.) Say it once more.
\end{tcolorbox}
\section[\texorpdfstring{\hi{दरोगा} (darogā)}{दरोगा (darogā)}]{\hi{दरोगा} (darogā) — a police inspector}

\label{lesson:HI-C290-daroga}



\begin{quote}
Before the new one: say the Hindi for a report (to the police), then the Hindi for to have written down.
\end{quote}

\subsection*{You'll want to know: \hi{दरोगा}}
\textbf{\hi{दरोगा}} — \emph{darogā} — \textquotedblleft{}a police inspector\textquotedblright{}.

\begin{quote}
\hi{दरोगा} \hi{ने} \hi{चौकीदार} \hi{से} \hi{बात} \hi{की।}
\end{quote}

\begin{quote}
\emph{darogā ne caukīdār se bāt kī.}
\end{quote}

\begin{quote}
The inspector talked to the watchman.
\end{quote}

Masculine.

\begin{cousinweb}
\textbf{\hi{दरोगा}} (\emph{darogā}) is usually traced to Persian \emph{dārōġa}, \textquotedblleft{}an overseer\textquotedblright{}, a title that once named many kinds of official. In Hindi it is the inspector in charge of a police station.
\end{cousinweb}

\begin{grammarlens}[title={What you've built}]
One more word for telling what happened.
\end{grammarlens}

\subsection*{Your turn}
\begin{itemize}
\raggedright
  \item \emph{Say it:} \emph{darogā}
  \item \emph{Say it:} \emph{darogā}, once more
  \item \emph{Say it:} \emph{darogā ne caukīdār se bāt kī.} — The inspector talked to the watchman.
  \item \emph{Recall:} say the Hindi for a report (to the police), then the Hindi for to have written down, then say \emph{darogā} again
\end{itemize}

\begin{tcolorbox}[breakable,colback=teal!4,colframe=teal!35!black,title={Before you move on}]
What does \hi{दरोगा} mean? (\textquotedblleft{}A police inspector\textquotedblright{}.) Say it once more.
\end{tcolorbox}
\section[\texorpdfstring{\hi{हवालात} (havālāt)}{हवालात (havālāt)}]{\hi{हवालात} (havālāt) — a lock-up}

\label{lesson:HI-C290-havalat}



\begin{quote}
Before the new one: say the Hindi for to have written down, then the Hindi for a police inspector.
\end{quote}

\subsection*{You'll want to know: \hi{हवालात}}
\textbf{\hi{हवालात}} — \emph{havālāt} — \textquotedblleft{}a lock-up\textquotedblright{}.

\begin{quote}
\hi{पुलिस} \hi{ने} \hi{चोर} \hi{को} \hi{रात} \hi{भर} \hi{हवालात} \hi{में} \hi{रखा।}
\end{quote}

\begin{quote}
\emph{pulis ne cor ko rāt bhar havālāt meṁ rakhā.}
\end{quote}

\begin{quote}
The police kept the thief in the lock-up all night.
\end{quote}

Feminine.

\begin{cousinweb}
is usually traced, through Persian, to Arabic \emph{ḥawāla}, \textquotedblleft{}a handing over, custody\textquotedblright{}: the room where a person handed over to the police is held. It is not \textbf{\hi{जेल}}, prison: \textbf{\hi{हवालात}} is for a night or two.
\end{cousinweb}

\begin{grammarlens}[title={What you've built}]
One more word for telling what happened.
\end{grammarlens}

\subsection*{Your turn}
\begin{itemize}
\raggedright
  \item \emph{Say it:} \emph{havālāt}
  \item \emph{Say it:} \emph{havālāt}, once more
  \item \emph{Say it:} \emph{pulis ne cor ko rāt bhar havālāt meṁ rakhā.} — The police kept the thief in the lock-up all night.
  \item \emph{Recall:} say the Hindi for to have written down, then the Hindi for a police inspector, then say \emph{havālāt} again
\end{itemize}

\begin{tcolorbox}[breakable,colback=teal!4,colframe=teal!35!black,title={Before you move on}]
What does \hi{हवालात} mean? (\textquotedblleft{}A lock-up\textquotedblright{}.) Say it once more.
\end{tcolorbox}
